\documentclass[11pt,a4paper]{article}

\usepackage[utf8]{inputenc}
\usepackage[T1]{fontenc}
\usepackage{lmodern}
\usepackage{amsmath,amssymb,amsthm}
\usepackage[margin=2.5cm]{geometry}
\usepackage[hidelinks]{hyperref}

\theoremstyle{definition}
\newtheorem{definition}{Definição}[section]

\theoremstyle{plain}
\newtheorem{theorem}[definition]{Teorema}
\newtheorem{lemma}[definition]{Lema}
\newtheorem{corollary}[definition]{Corolário}

\title{Demonstração de Lea Overleaf}
\author{Ximena Gomez}
\date{Outubro de 2026}

\begin{document}

\maketitle

\section*{Objetivo da demonstração}

Vamos usar Lea Overleaf para transformar definições e enunciados
matemáticos em código Lean 4.

A demonstração começa com uma função simples e avança para provas
que reutilizam resultados anteriores. Em uma delas, pedimos
explicitamente que a prova seja dividida em etapas locais.

Os exemplos de \texttt{possibilities}, \texttt{Day} e
\texttt{isWeekend} foram adaptados de
\href{https://emap-nlp.github.io/book/}
{Semântica computacional com Lean}.

\section{Contagem de possibilidades}

Se cada fato básico pode ser verdadeiro ou falso, cada novo fato
duplica a quantidade de possibilidades.

\subsection{Criar uma definição}

\begin{definition}[Função de contagem]
\label{possibilities}
A função recursiva chamada \texttt{possibilities}, tal que
\[
  \operatorname{possibilities}(0)=1
\]
e, para todo número natural $n$,
\[
  \operatorname{possibilities}(n+1)
  =2\operatorname{possibilities}(n).
\]
\end{definition}


\subsection{Provar uma propriedade imediata}

\begin{lemma}[Caso base]
\label{possibilities_zero}
A função de contagem satisfaz
\[
  \operatorname{possibilities}(0)=1.
\]
\end{lemma}

\subsection{Provar por indução}

\begin{theorem}[Forma fechada]
\label{possibilities_pow}

Para todo número natural $n$,
\[
  \operatorname{possibilities}(n)=2^n.
\]
\end{theorem}

\subsection{Reutilizar um teorema}

\begin{corollary}[Dez fatos]
\label{possibilities_ten}
Para dez fatos básicos, existem $1024$ possibilidades:
\[
  \operatorname{possibilities}(10)=1024.
\]
\end{corollary}

\begin{lemma}[Positividade]
\label{possibilities_positive}\
Para todo número natural $n$,
\[
  0<\operatorname{possibilities}(n).
\]
\end{lemma}

\subsection{Dividir uma prova em etapas}

\begin{theorem}[Combinação de grupos de fatos]
\label{possibilities_add}
Para todos os números naturais $n$ e $m$,
\[
  \operatorname{possibilities}(n+m)
  =\operatorname{possibilities}(n)
   \operatorname{possibilities}(m).
\]
\end{theorem}


\begin{corollary}[Dois grupos do mesmo tamanho]
\label{possibilities_double}
% lea: formalize label=possibilities_double uses={possibilities, possibilities_add} context={Use possibilities_add com os dois argumentos iguais a n. Relacione o produto de um numero por si mesmo com seu quadrado. Nao refaca a prova por inducao.}
Para todo número natural $n$,
\[
  \operatorname{possibilities}(n+n)
  =\bigl(\operatorname{possibilities}(n)\bigr)^2.
\]
\end{corollary}

\section{Dias da semana}

Agora vamos criar um tipo indutivo e uma função booleana.

\begin{definition}[Dias e fim de semana]
\label{isWeekend}
% lea: define label=isWeekend context={No mesmo arquivo, crie o tipo Day e depois a funcao isWeekend. Use casamento de padroes. Mantenha isWeekend como declaracao principal.}
Defina um tipo indutivo chamado \texttt{Day}, com exatamente
os sete construtores:
\[
\begin{gathered}
\texttt{monday},\quad \texttt{tuesday},\quad
\texttt{wednesday},\quad \texttt{thursday},\\
\texttt{friday},\quad \texttt{saturday},\quad
\texttt{sunday}.
\end{gathered}
\]

Defina também a função \texttt{isWeekend},
de \texttt{Day} para \texttt{Bool}.
Ela deve retornar \texttt{true} para sábado e domingo,
e \texttt{false} para os demais dias.
\end{definition}

\begin{theorem}[Caracterização do fim de semana]
\label{isWeekend_iff}
Para todo $d$ do tipo \texttt{Day},
\[
  \texttt{isWeekend}\ d=\texttt{true}
  \quad\Longleftrightarrow\quad
  d=\texttt{Day.saturday}
  \ \lor\
  d=\texttt{Day.sunday}.
\]
\end{theorem}

\begin{theorem}[Dias que não são fim de semana]
\label{notWeekend_iff}
Para todo $d$ do tipo \texttt{Day},
\[
\begin{aligned}
  \texttt{isWeekend}\ d=\texttt{false}
  \quad\Longleftrightarrow\quad
  &d\ne\texttt{Day.saturday}\\
  &{}\land d\ne\texttt{Day.sunday}.
\end{aligned}
\]
\end{theorem}

\section{ uma afirmação falsa}

\begin{theorem}[Afirmação ]
\label{every_day_weekend}
Para todo dia $d$ do tipo \texttt{Day},
\[
  \texttt{isWeekend}\ d=\texttt{true}.
\]
\end{theorem}

\end{document}